%=====================================================================
%  CityVision — Hugging Face Spaces deployment guide (PDF version)
%  Compile:  pdflatex DEPLOYMENT.tex   (run once; no cross-refs)
%=====================================================================
\documentclass[11pt]{article}

\usepackage[utf8]{inputenc}
\usepackage[T1]{fontenc}
\usepackage[a4paper,margin=2.2cm]{geometry}
\usepackage{lmodern}
\usepackage{microtype}
\usepackage{xcolor}
\usepackage{booktabs}
\usepackage{tabularx}
\usepackage{enumitem}
\usepackage{listings}
\usepackage{amssymb}
\usepackage{titlesec}
\usepackage[colorlinks=true,urlcolor=blue,linkcolor=black]{hyperref}

\definecolor{cvblue}{RGB}{40,90,150}
\definecolor{codebg}{RGB}{244,246,248}
\definecolor{rulegray}{RGB}{120,120,120}

% ---- code styling ----
\lstset{
  basicstyle=\ttfamily\footnotesize,
  breaklines=true,
  breakatwhitespace=false,
  columns=fullflexible,
  keepspaces=true,
  showstringspaces=false,
  upquote=true,
}
\lstdefinestyle{block}{
  backgroundcolor=\color{codebg},
  frame=single,
  rulecolor=\color{rulegray!40},
  framesep=5pt,
  xleftmargin=6pt,
  xrightmargin=2pt,
  basicstyle=\ttfamily\small,
}
% inline code shortcut:  \c{some_code}
\newcommand{\code}[1]{\lstinline[basicstyle=\ttfamily\normalsize]|#1|}

% ---- headings ----
\titleformat{\section}{\Large\bfseries\color{cvblue}}{\thesection.}{0.5em}{}
\titleformat{\subsection}{\large\bfseries}{\thesubsection}{0.5em}{}
\setlist[itemize]{leftmargin=1.4em,itemsep=2pt,topsep=2pt}
\setlist[enumerate]{leftmargin=1.6em,itemsep=2pt,topsep=2pt}

\newcommand{\rarr}{$\rightarrow$}

\title{\vspace{-1.2cm}\textbf{\color{cvblue}Deploying CityVision on Hugging Face Spaces}\\[2pt]
\large A complete, repo-accurate guide}
\author{CityVision — urban scene semantic segmentation}
\date{\today}

\begin{document}
\maketitle
\vspace{-0.6cm}

\noindent
This guide deploys CityVision as \textbf{two Hugging Face Spaces}:
\begin{itemize}
  \item \textbf{API Space} — the FastAPI prediction service (carries the model
    weights, uses PyTorch).
  \item \textbf{App Space} — the Streamlit demo UI (lean, no PyTorch; calls the API
    over HTTP).
\end{itemize}
It covers the \emph{why} (separation, Docker) and every concrete step: creating
Spaces, git-lfs for the weights, what to push where, the port gotcha, wiring the
two together, updating, and troubleshooting.

\medskip
\noindent\textbf{Repo references:} \code{deploy/hf/api/Dockerfile},
\code{deploy/hf/api/README.md}, \code{deploy/hf/app/Dockerfile},
\code{deploy/hf/app/README.md}, \code{docker/*.Dockerfile},
\code{docker-compose.yml}, \code{.dockerignore}.

%=====================================================================
\section{Mental model — how HF Spaces work}

A Hugging Face \textbf{Space is just a git repository} hosted on
\code{huggingface.co}, with a few extra rules:

\begin{enumerate}
  \item \textbf{A \code{README.md} with a YAML front-matter header configures the
    Space.} The key fields:
\begin{lstlisting}[style=block]
---
title: CityVision Segmentation API
emoji: <road emoji>
colorFrom: blue
colorTo: green
sdk: docker        # build a Dockerfile (not the Gradio/Streamlit auto-SDK)
app_port: 7860     # the port HF routes public traffic to
pinned: false
---
\end{lstlisting}
  \item \textbf{\code{sdk: docker} means HF builds the \code{Dockerfile} at the repo
    root} and runs the container. Whatever your Dockerfile's \code{CMD} starts is
    your app.
  \item \textbf{HF exposes exactly one port: \code{app_port} (7860 by convention).}
    Your server \emph{must} listen on it. This is why the HF Dockerfiles serve on
    7860 while local \code{docker-compose} uses 8000/8501 (see the port gotcha).
  \item \textbf{Public URLs:}
    \begin{itemize}
      \item Space page (build logs, settings):
        \code{https://huggingface.co/spaces/<user>/<space>}
      \item Direct app URL (what clients hit): \code{https://<user>-<space>.hf.space}
        (username + space, lowercased, non-alphanumerics \rarr{} \code{-}).
      \item Example: user \code{DaryaEL}, space \code{cityvision-api} \rarr{}
        \code{https://daryael-cityvision-api.hf.space}.
    \end{itemize}
  \item \textbf{Files $>$ 10 MB must be tracked with git-lfs} (our ResNet50 weight is
    $\sim$157 MB).
\end{enumerate}

You end up with \textbf{two independent Space repos}, each with a root-level
\code{Dockerfile} + \code{README.md} and the subset of the project it needs.

%=====================================================================
\section{Why separate frontend and backend?}

\begin{tabularx}{\linewidth}{@{}p{3.2cm}X@{}}
\toprule
\textbf{Reason} & \textbf{Detail} \\
\midrule
Different dependency weight & Backend needs PyTorch + torchvision ($\sim$hundreds of
MB) and the weights; frontend needs only Streamlit + requests + Pillow + numpy ---
\emph{no torch}. Splitting keeps the UI image small and fast. \\
\addlinespace
Independent scaling \& cost & Inference is the expensive part. Put the \textbf{API on
paid always-on} hardware (no cold starts) while the \textbf{UI stays free}. A
monolith couldn't. \\
\addlinespace
Clean contract & The frontend only makes HTTP calls (\code{/predict/image},
\code{/predict/classes}, \code{/health}). Any client --- curl, browser, another app
--- uses the API identically. \\
\addlinespace
No training/serving skew & API and training import the \textbf{same
\code{cityvision/} package} (classes, palette, model defs), so preprocessing is
identical everywhere. \\
\bottomrule
\end{tabularx}

\medskip
The frontend reaches the backend purely through the \code{CITYVISION_API}
environment variable --- locally \code{http://backend:8000} (compose network), in the
cloud the API Space's public URL.

%=====================================================================
\section{Why Docker (and why the Docker SDK on HF)?}

\begin{itemize}
  \item \textbf{Reproducibility} --- the image pins exact versions
    (\code{torch==2.12.0}, \code{fastapi==0.136.3}, \dots). "Works on my machine"
    becomes "works everywhere".
  \item \textbf{Same artifact locally and in the cloud} --- validate with
    \code{docker compose up}, then deploy the same Dockerfile logic to HF.
  \item \textbf{Full control of the environment} --- we install \textbf{CPU-only
    PyTorch wheels} (\code{--index-url https://download.pytorch.org/whl/cpu}), far
    smaller than the CUDA build. HF free/basic hardware is CPU anyway.
  \item \textbf{Portability / no lock-in} --- the same images run on
    AWS/GCP/Azure/Render. HF is a convenient host, not a dependency.
  \item \textbf{HF \code{sdk: docker}} gives this control instead of HF's opinionated
    auto-SDKs --- we choose the base image, deps, and start command.
\end{itemize}

%=====================================================================
\section{Prerequisites (one-time)}

\begin{enumerate}
  \item A \textbf{Hugging Face account} --- \url{https://huggingface.co/join}
  \item \textbf{git} and \textbf{git-lfs} installed:
\begin{lstlisting}[style=block]
git lfs version    # should print a version; install git-lfs if not
git lfs install    # one-time, enables the lfs filters for your user
\end{lstlisting}
  \item A \textbf{trained weight} at \code{backend/model/resnet50_best.pth} (the
    served model; \code{CITYVISION_ARCH=ResNet50-UNet}).
  \item An \textbf{HF access token} with \emph{write} permission (HF \rarr{} Settings
    \rarr{} Access Tokens \rarr{} New token, role \textbf{Write}). Use it as the git
    password when pushing (username = your HF username).
  \item \emph{(Optional)} \textbf{Docker Desktop} to test locally first.
\end{enumerate}

%=====================================================================
\section{Step 0 — Sanity-check locally with docker compose}

Before touching the cloud, confirm the two-service setup works end to end:
\begin{lstlisting}[style=block]
docker compose up --build
# open http://localhost:8501  (UI)   and   http://localhost:8000/docs  (API)
\end{lstlisting}

This proves the backend image builds and loads
\code{backend/model/resnet50_best.pth}; the frontend image builds (no torch) and
talks to the backend via \code{CITYVISION_API=http://backend:8000} (set in
\code{docker-compose.yml}). Stop with \code{docker compose down}.

%=====================================================================
\section{Step 1 — Deploy the API Space}

\subsection{Create the Space (web UI)}
\begin{enumerate}
  \item \url{https://huggingface.co/new-space}
  \item \textbf{Owner}: you. \textbf{Space name}: \code{cityvision-api}.
  \item \textbf{SDK}: \textbf{Docker} \rarr{} Blank/from scratch (not a template).
  \item \textbf{Hardware}: start with \textbf{CPU basic (free)}; upgrade to paid
    always-on later to avoid cold starts.
  \item \textbf{Visibility}: Public. Create --- HF makes an (almost) empty git repo.
\end{enumerate}

\subsection{Populate the Space repo}
The API Space needs, \textbf{at its repo root}:
\begin{lstlisting}[style=block]
Dockerfile          <- from deploy/hf/api/Dockerfile
README.md           <- from deploy/hf/api/README.md (has the HF front-matter)
cityvision/         <- shared package (constants + model defs)
backend/            <- FastAPI app + backend/model/resnet50_best.pth (via LFS!)
\end{lstlisting}
\emph{Why these:} the API Dockerfile does \code{COPY cityvision/} and
\code{COPY backend/}, sets \code{ENV CITYVISION_ARCH=ResNet50-UNet}, and serves
\code{uvicorn backend.app:app} on 7860. It does \emph{not} need \code{frontend/},
\code{data/}, \code{notebooks/}, \code{mlflow.db}.

\begin{lstlisting}[style=block]
# 1. Clone the empty Space repo (use your HF token as the password)
git clone https://huggingface.co/spaces/<user>/cityvision-api
cd cityvision-api

# 2. Enable LFS and TRACK THE WEIGHTS *before* adding them
git lfs install
git lfs track "*.pth"          # creates/updates .gitattributes

# 3. Copy the needed files from your project (adjust the source path)
cp   /path/to/CityVision/deploy/hf/api/Dockerfile ./Dockerfile
cp   /path/to/CityVision/deploy/hf/api/README.md  ./README.md
cp -r /path/to/CityVision/cityvision ./cityvision
cp -r /path/to/CityVision/backend    ./backend
ls -lh backend/model/resnet50_best.pth   # ~157 MB

# 4. Commit and push
git add .gitattributes Dockerfile README.md cityvision backend
git commit -m "Deploy CityVision prediction API (Docker Space)"
git push
\end{lstlisting}

\noindent\textbf{Critical ordering:} run \code{git lfs track "*.pth"}
\emph{before} \code{git add}-ing the weight. Adding a $>$10 MB file without LFS gets
the push rejected. Verify with \code{git lfs ls-files} --- the \code{.pth} must be
listed.

\subsection{Watch the build \& test}
Open the Space's \textbf{Logs} tab (you'll see \code{pip install torch}, the
\code{COPY}s, then \code{uvicorn ... 0.0.0.0:7860}). Then:
\begin{lstlisting}[style=block]
curl https://<user>-cityvision-api.hf.space/health
# {"status":"ok","arch":"ResNet50-UNet", ...}

# Swagger UI:  https://<user>-cityvision-api.hf.space/docs

curl -X POST "https://<user>-cityvision-api.hf.space/predict/image?format=overlay" \
     -F "file=@some_street.png" -o overlay.png
\end{lstlisting}
\textbf{Note this URL} --- the App Space needs it.

%=====================================================================
\section{The port gotcha (read this)}

\begin{tabularx}{\linewidth}{@{}Xccp{4.2cm}@{}}
\toprule
\textbf{Environment} & \textbf{Backend} & \textbf{Frontend} & \textbf{Set where} \\
\midrule
Local \code{docker-compose} & 8000 & 8501 & \code{docker/*.Dockerfile},
\code{docker-compose.yml} \\
\addlinespace
\textbf{Hugging Face Space} & \textbf{7860} & \textbf{7860} &
\code{deploy/hf/*/Dockerfile} + \code{app_port: 7860} in each \code{README.md} \\
\bottomrule
\end{tabularx}

\medskip
HF routes public traffic only to \code{app_port} (7860). Each container serves a
\textbf{single} service on 7860 --- the whole reason we run two Spaces. The
\code{EXPOSE 7860} and \code{--port 7860} in the HF Dockerfiles must match
\code{app_port} in the front-matter.

%=====================================================================
\section{Step 2 — Deploy the App Space}

\subsection{Create it}
Same as before, name it \code{cityvision}. SDK = \textbf{Docker}, CPU basic (free) is
fine.

\subsection{Populate \& point it at the API}
At the repo root:
\begin{lstlisting}[style=block]
Dockerfile          <- from deploy/hf/app/Dockerfile
README.md           <- from deploy/hf/app/README.md
cityvision/         <- torch-free constants (palette/classes)
frontend/           <- Streamlit UI
data/samples/       <- committed sample images (demo works w/o dataset)
\end{lstlisting}
\begin{lstlisting}[style=block]
git clone https://huggingface.co/spaces/<user>/cityvision
cd cityvision
git lfs install

cp   /path/to/CityVision/deploy/hf/app/Dockerfile ./Dockerfile
cp   /path/to/CityVision/deploy/hf/app/README.md  ./README.md
cp -r /path/to/CityVision/cityvision   ./cityvision
cp -r /path/to/CityVision/frontend     ./frontend
mkdir -p data && cp -r /path/to/CityVision/data/samples ./data/samples

git add .
git commit -m "Deploy CityVision demo UI (Docker Space)"
git push
\end{lstlisting}

\subsection{Set the API URL}
The app finds the backend via \code{CITYVISION_API}. Two ways:
\begin{itemize}
  \item \textbf{Simplest:} edit the App Space's Dockerfile line
    \code{ENV CITYVISION_API=https://daryael-cityvision-api.hf.space} to your API
    URL, commit/push. (Already the default in \code{deploy/hf/app/Dockerfile}.)
  \item \textbf{Cleaner:} Space \rarr{} Settings \rarr{} Variables and secrets \rarr{}
    add a \textbf{Variable} \code{CITYVISION_API} =
    \code{https://<user>-cityvision-api.hf.space}. Space env vars override the
    Dockerfile \code{ENV}.
\end{itemize}
Open \code{https://huggingface.co/spaces/<user>/cityvision}, pick an image \rarr{} it
calls your API and shows \textbf{real image \textperiodcentered{} real mask
\textperiodcentered{} predicted mask}.

\medskip
\noindent The frontend tolerates a sleeping API: \code{wait_for_api()} polls
\code{/health} through the $\sim$30--60 s cold-start window (503s) instead of
failing.

%=====================================================================
\section{git-lfs — the model weight, in detail}

\begin{itemize}
  \item HF rejects files $>$ 10 MB unless tracked by LFS. \code{resnet50_best.pth} is
    $\sim$157 MB.
  \item \code{git lfs track "*.pth"} writes the \code{.gitattributes} line
    \code{*.pth filter=lfs diff=lfs merge=lfs -text} that routes the weight through
    LFS. \textbf{Commit \code{.gitattributes} in the same commit as the weight.}
  \item Verify before pushing: \code{git lfs ls-files} must show the \code{.pth}.
  \item If you committed the weight \emph{without} LFS: fastest fix is delete the
    local clone, re-clone, \code{git lfs track} first, then re-copy and re-commit
    (avoids history surgery).
  \item The full Cityscapes dataset is \emph{not} deployed (gitignored + excluded in
    \code{.dockerignore}); only \code{data/samples/} ships in the App image.
\end{itemize}

%=====================================================================
\section{Hardware, cold starts, and "always-on"}

\begin{itemize}
  \item \textbf{Free CPU basic} Spaces \textbf{sleep after inactivity} and take
    $\sim$30--60 s to wake ("cold start"), returning 503 meanwhile.
  \item For a smooth demo, run the \textbf{API on paid always-on} CPU so the model
    stays resident. The \textbf{App} can stay free (cheap cold start; it waits for
    the API via \code{wait_for_api}).
  \item Change in Space \rarr{} Settings \rarr{} Hardware. Keep the two Spaces
    separate --- this is the flexibility the split buys you.
\end{itemize}

%=====================================================================
\section{Configuration reference (env vars)}

\begin{tabularx}{\linewidth}{@{}llXp{3.3cm}@{}}
\toprule
\textbf{Variable} & \textbf{Used by} & \textbf{Meaning} & \textbf{Default} \\
\midrule
\code{CITYVISION_API} & App & URL of the API Space & \code{http://backend:8000}
(local) / your API URL (cloud) \\
\addlinespace
\code{CITYVISION_ARCH} & API & Which architecture to serve & \code{ResNet50-UNet} \\
\addlinespace
\code{CITYVISION_CHECKPOINT} & API & Explicit weight path (optional) & derived from
arch \\
\addlinespace
\code{CITYVISION_MAX_UPLOAD_MB} & API & Upload size cap & \code{10} \\
\bottomrule
\end{tabularx}

\medskip
Set these under Space \rarr{} Settings \rarr{} Variables and secrets. Use a
\emph{secret} only for truly sensitive values --- none here are secret.

%=====================================================================
\section{Updating a deployment}

HF redeploys on every push:
\begin{lstlisting}[style=block]
cd cityvision-api        # or cityvision
# ...make changes / drop in a newly trained weight...
git add -A && git commit -m "Update model / code" && git push
# HF rebuilds the image and restarts the container automatically.
\end{lstlisting}
To ship a \textbf{new weight}: replace \code{backend/model/resnet50_best.pth} in the
API Space repo, commit (LFS handles it), push. No code change needed.

%=====================================================================
\section{Troubleshooting}

\begin{tabularx}{\linewidth}{@{}XX@{}}
\toprule
\textbf{Symptom} & \textbf{Likely cause / fix} \\
\midrule
Build fails on push: "file too large" & Weight not tracked by LFS.
\code{git lfs track "*.pth"} \emph{before} adding; check \code{git lfs ls-files}. \\
\addlinespace
Builds but "no application on port" & Server not on \textbf{7860}. Ensure
\code{--port 7860} in \code{CMD} and \code{app_port: 7860} in \code{README.md}. \\
\addlinespace
API \code{/health} OK but App can't connect & \code{CITYVISION_API} wrong/missing.
Set it to \code{https://<user>-<space>.hf.space} (no trailing slash, https). \\
\addlinespace
Predictions 503 for a while then recover & API cold start. Expected on free tier;
\code{wait_for_api} retries. Use always-on to avoid. \\
\addlinespace
\code{ModuleNotFoundError: cityvision} & \code{cityvision/} not copied into the Space
repo root. \\
\addlinespace
Model file missing at runtime & \code{backend/model/*.pth} didn't ship (LFS pointer
only, or not copied). Confirm real bytes. \\
\addlinespace
Wrong model served & \code{CITYVISION_ARCH} mismatch with the weight present. \\
\bottomrule
\end{tabularx}

%=====================================================================
\section{Deployment checklist}

\noindent\textbf{API Space}
\begin{itemize}[label=$\square$]
  \item Space created, SDK = Docker, port 7860
  \item \code{Dockerfile} (from \code{deploy/hf/api/}) + \code{README.md} at root
  \item \code{cityvision/} and \code{backend/} copied in
  \item \code{git lfs track "*.pth"} done \emph{before} adding the weight;
    \code{.gitattributes} committed
  \item \code{backend/model/resnet50_best.pth} present as real bytes
    (\code{git lfs ls-files})
  \item Build green; \code{/health} and \code{/docs} reachable
  \item (optional) upgraded to always-on hardware
\end{itemize}

\noindent\textbf{App Space}
\begin{itemize}[label=$\square$]
  \item Space created, SDK = Docker, port 7860
  \item \code{Dockerfile} (from \code{deploy/hf/app/}) + \code{README.md} at root
  \item \code{cityvision/}, \code{frontend/}, \code{data/samples/} copied in
  \item \code{CITYVISION_API} set to the API Space URL (Dockerfile ENV or Space
    variable)
  \item UI loads and shows image \textperiodcentered{} real mask
    \textperiodcentered{} predicted mask
\end{itemize}

%=====================================================================
\section*{TL;DR}
Two Docker Spaces. Each is a git repo with a root \code{Dockerfile} +
\code{README.md} (front-matter, \code{app_port: 7860}). API Space carries
\code{cityvision/} + \code{backend/} + the LFS-tracked weight and serves FastAPI on
7860. App Space carries \code{cityvision/} + \code{frontend/} + \code{data/samples/},
serves Streamlit on 7860, and points at the API via \code{CITYVISION_API}. Test
locally with \code{docker compose up}, then \code{git push} each Space.

\vfill
\noindent\footnotesize\color{rulegray}
Source: \url{https://github.com/delnouty/cityvision-segmentation} \quad\textbar\quad
Markdown version: \code{deploy/hf/DEPLOYMENT.md}

\end{document}
